\section{Percepción de datos y panel de evaluación}
\subsection{Vistas centrales del panel}
Graham recomienda construir un conjunto de paneles de evaluación de Agentes que monitoreen en tiempo real la tasa de éxito, el número de intervenciones, la duración media de ejecución y los registros de activación de las compuertas de seguridad. Cada indicador del panel debería vincularse con el log del evento correspondiente, para poder rastrear la causa raíz cuando un indicador se comporte de manera anómala.

\begin{table}[H]
\centering
\begin{tabular}{p{0.28\textwidth} p{0.42\textwidth} p{0.28\textwidth}}
\toprule
Indicador & Descripción & Objetivo \\
\midrule
Tasa de aprobación & Proporción en que el Agente ejecuta un patch y pasa las pruebas & \(\geq 40\%\) (Lite), \(\geq 20\%\) (Full) \\
\midrule
Tasa de intervención humana & Proporción que requiere pausa o aprobación de una persona & \(\leq 10\%\) \\
\midrule
Duración de ejecución & Tiempo medio desde Plan hasta la finalización de Verify & \(\leq 20 \text{ minutos}\) \\
\midrule
Activación de seguridad & Número de comandos que involucran la base de datos o el despliegue & Requiere revisión humana \\
\bottomrule
\end{tabular}
\caption{Indicadores centrales que se vigilan en el panel del Agente.}
\end{table}

\begin{knowledgebox}{Dependencias de datos del panel}
Cada indicador debería vincularse con el log del evento (comandos, diff, salida de las pruebas), para facilitar el rastreo y entrenar modelos interpretables.
\end{knowledgebox}

\subsection{Ciclo de retroalimentación del panel}
El panel no solo debe monitorear, también debe sustentar un proceso de «retroalimentación -> ajuste»: cuando un indicador retrocede, el sistema envía automáticamente el log de la falla al equipo de Agentes, para que produzcan una nueva combinación de prompt y herramientas, y la nueva estrategia se despliegue al entorno de pruebas.

\begin{importantbox}{Práctica de retroalimentación de datos}
Diseña reportes automatizados: cuando un indicador sea anómalo (por ejemplo, tasa de aprobación < 30\%), abre automáticamente un issue, adjunta el diff y el log, y notifica al equipo, para evitar que el problema quede olvidado por mucho tiempo.
\end{importantbox}

\subsection{Resumen de la sección}
\begin{itemize}
    \item El panel en tiempo real es una ventana clave para juzgar la madurez de un Agente.
    \item Convertir los eventos anómalos en un ciclo de retroalimentación automatizado permite acelerar la iteración de la estrategia.
\end{itemize}

